\textbf{Время:} резерв

\textbf{Тезис:} архивные зависимости Rout и Rside показывают, как параметры поперечной области испускания меняются с импульсом пары и центральностью.

\textbf{На экране:} названия панелей — Rout слева, Rside справа; общие оси kT и радиус в фм; нижняя цветовая легенда центральностей; проследить красные точки слева направо в каждой панели, затем сопоставить цвета при близком kT.

\textbf{Текст:} От корреляционных проекций переходим к параметрам их ширины. Слева показан радиус out, справа — side. По горизонтали на обеих панелях идёт k-T, по вертикали — радиус в фемтометрах.

Сначала прочитаем нижнюю легенду. Красный цвет обозначает центральность от нуля до десяти процентов, синий — от десяти до тридцати, розовый — от тридцати до пятидесяти, зелёный — от пятидесяти до восьмидесяти. Здесь цвет различает группы событий; в предыдущем объединённом графике STAR он различал направления.

Проследим красные точки слева направо. В этой архивной обработке оба поперечных радиуса с ростом k-T уменьшаются. Затем сравним разные цвета при близком k-T: для разных центральностей получены разные масштабы.

Теперь сопоставим левую и правую панели. Out и side дают разные численные значения. Это согласуется с тем, что они описывают разные характеристики выбранной области испускания: для out существенна также связь положения и времени. Поэтому эти два параметра сохраняют по отдельности.

Перед нами результаты прежней обработки модельных событий. Для прямого сравнения с экспериментом ещё потребуется согласовать условия анализа.

\textbf{Переход:} Посмотрим, как в той же работе выглядит третий радиус — вдоль пучка.

\textbf{Справка:} Балинов (2025), исходный слайд 17, две верхние панели. Оригинальные точки и погрешности сохранены; их новая оценка не выполнялась. Гауссовы радиусы характеризуют области однородности выбранных пар. Различие Rout и Rside само по себе не даёт длительность испускания: нужны пространственные дисперсии и пространственно-временная ковариация. Научный смысл соответствующих моментов разобран у Lisa et al. (2005), §2.5.2. В архиве не установлены код фита, полные настройки, метод погрешностей, заряд именно этих панелей и способ калибровки центральности; численный масштаб старого фита независимо не подтверждён.
